\section{Cadenas singulares lineales}

% ------------------------------------------------------------------------
\subsection{Símplices singulares lineales}

\begin{frame}
  \frametitle{Símplices singulares lineales}

  En toda esta parte, $Y\subseteq\R^N$ es un subconjunto \alert{convexo}.

  \begin{definition}
    Un \alert{símplice singular lineal} de dimensión $n$ en $Y$ es una
    aplicación $\sigma\colon\simplex{n}\to Y$ que es la restricción de una
    aplicación lineal $\R^{n+1}\to\R^N$.
  \end{definition}

  Por linealidad, $\sigma$ está determinado por los vértices
  $v_i=\sigma(e_i)\in Y$, y
  \[
    \sigma(t_0,\ldots,t_n)=t_0v_0+\cdots+t_nv_n.
  \]
  Recíprocamente, como $Y$ es convexo, cualesquiera $v_0,\ldots,v_n\in Y$
  definen así un símplice singular lineal. Los símplices singulares lineales
  en $Y$ son \alert{lo mismo} que las listas de $n+1$ puntos de $Y$, y su
  imagen es el símplice generalizado $[v_0,\ldots,v_n]\subseteq Y$. Por eso
  escribimos
  \[
    \sigma=[v_0,\ldots,v_n].
  \]
\end{frame}

\begin{frame}
  \frametitle{Ejemplos}

  \begin{itemize}
    \item La identidad de $\simplex{n}$ es el símplice singular lineal
      \[
        \id_{\simplex{n}}=[e_0,\ldots,e_n]\colon
        \simplex{n}\longrightarrow\simplex{n}.
      \]
    \item Las inclusiones de las caras
      $\delta_i\colon\simplex{n-1}\to\simplex{n}$ son
      $\delta_i=[e_0,\ldots,\widehat{e}_i,\ldots,e_n]$.
    \item Si $\sigma=[v_0,\ldots,v_n]$, la cara $\sigma\delta_i$ es el
      símplice singular lineal $\partial_i\sigma=[v_0,\ldots,\widehat{v}_i,\ldots,v_n]$.
  \end{itemize}
\end{frame}

% ------------------------------------------------------------------------
\subsection{El complejo de cadenas lineales}

\begin{frame}
  \frametitle{Cadenas singulares lineales}

  \begin{definition}
    El grupo de \alert{cadenas singulares lineales}
    $\LC_n(Y)\subseteq\Sing_n(Y)$ es el subgrupo generado por los
    símplices singulares lineales de dimensión $n$ en $Y$.
  \end{definition}

  \medskip

  Es abeliano libre, con base los símplices singulares lineales en $Y$.
  Las caras de un símplice singular lineal son símplices singulares
  lineales, luego el borde de una cadena singular lineal es una cadena singular
  lineal: $\LC_\bullet(Y)$ es un subcomplejo de $\Sing_\bullet(Y)$, el
  \alert{complejo de cadenas singulares lineales} de $Y$.
\end{frame}

% ------------------------------------------------------------------------
\subsection{La versión aumentada}

\begin{frame}
  \frametitle{Cadenas singulares lineales aumentadas}

  \begin{definition}
    El complejo \alert{aumentado} $\widetilde{\LC}_\bullet(Y)$ coincide con
    $\LC_\bullet(Y)$ en dimensiones $n\geq0$, y
    $\widetilde{\LC}_{-1}(Y)=\Z\,[\;]$, con
    $\partial[v_0]=[\;]$ para todo $v_0\in Y$.
  \end{definition}

  \medskip

  Es un subcomplejo del complejo de cadenas singulares aumentado
  $\widetilde{\Sing}_\bullet(Y)$. Convenimos que $\partial=0$ en grado $-1$.
\end{frame}
